\documentclass[11pt,a4paper]{article}
\usepackage[margin=2.5cm]{geometry}
\usepackage{amsmath,amssymb,amsthm}
\usepackage{xcolor}
\usepackage{booktabs}
\usepackage{enumitem}
\usepackage{hyperref}
\hypersetup{colorlinks=true, linkcolor=blue!60!black, urlcolor=blue!60!black, citecolor=blue!60!black}

\usepackage{titlesec}
\titleformat{\section}{\large\bfseries}{\thesection.}{0.5em}{}
\titleformat{\subsection}{\normalsize\bfseries}{\thesubsection.}{0.5em}{}

\setlength{\parindent}{0pt}
\setlength{\parskip}{6pt}

\title{\textbf{The Laplace--Beltrami Operator on Hurwitz\\
and Hurwitz-like Surfaces}\\[4pt]
\large Spectral Gaps, Representation Theory, and Topological Constraints}
\author{Robin Langer}
\date{}

\begin{document}
\maketitle

\begin{abstract}
The Laplace--Beltrami operator on a compact hyperbolic Riemann surface encodes
three things simultaneously: the geometry (curvature, metric), the topology
(genus, handles), and the symmetry (automorphism group). Its eigenvalues are
the natural frequencies of the surface; its eigenspaces are irreducible
representations of the symmetry group; and its first cohomology measures the
topological obstructions that no amount of geometric deformation can remove.
We survey this structure on a family of hyperbolic surfaces arising as quotients
of the hyperbolic plane by triangle groups---the Bolza surface (genus~2,
$\Delta(2,3,8)$, $\mathrm{GL}(2,\mathbb{F}_3)$), the Klein quartic (genus~3,
$\Delta(2,3,7)$, $\mathrm{PSL}(2,7)$), and the Macbeath surface (genus~7,
$\Delta(2,3,7)$, $\mathrm{PSL}(2,8)$)---and discuss the tension between
maximising symmetry and maximising spectral gap. Companion code computing the
cell complexes, homology, and group representations is included.
\end{abstract}

\section{Motivation: From LLMs to Hyperbolic Surfaces}

This paper is a companion to \emph{LLMs as Resonant Structures}~\cite{langer2026llm},
which develops a spectral theory of large language models in three layers:
Green's functions (linear spectral decomposition), the Inverse Scattering
Transform (nonlinear spectral decomposition via sparse autoencoders), and
sheaf theory (the global obstruction to interpretability).

The sheaf-theoretic layer raises a concrete mathematical question. If one treats
the LLM's latent space as a Riemannian manifold $M$, then local linearisations
(SAE decompositions that work in a neighbourhood of input space) form a
\emph{presheaf} on $M$. Polysemanticity---the phenomenon where a single feature
means different things in different contexts---is the failure of the gluing
axiom: the presheaf is not a sheaf. The first \v{C}ech cohomology $H^1(M, \mathcal{F})$
measures the degree of this failure.

The natural next question is: \emph{how much of $H^1$ is topologically forced?}
If the manifold $M$ has nontrivial topology (handles, genus), then Gauss--Bonnet
and Hodge theory guarantee that $H^1 \neq 0$ regardless of the metric---no
training objective can eliminate it. The Laplace--Beltrami operator
$\Delta_g$ on $M$ is the bridge between the topology ($H^1$), the geometry
(curvature, holonomy), and the representation theory (how the symmetry group
organises the eigenspaces).

I should be honest about what I don't know. My background is in symmetric graphs
and symmetric polynomials---algebraic combinatorics, not differential geometry.
I view everything through the lens of symmetry. I do not have working
intuition for the Laplace--Beltrami operator, and this paper is my attempt to
build some.

The surfaces studied here---the Bolza surface, the Klein quartic, the Macbeath
surface---were not chosen because they are the right model for an LLM's latent
manifold. They were chosen because I already know how to compute with them.
In prior work on quantum error-correcting codes~\cite{langer2026css}, I
formalised in Lean~4 the fact that a CSS code constructed from a surface tiling
encodes exactly $2g$ logical qubits, where $g$ is the genus. The chain complex
condition $\partial^2 = 0$ (which gives the orthogonality of X and Z
stabilizers) and the first Betti number $\beta_1 = 2g$ (which counts the
logical qubits) are the same $H^1$ that appears in the sheaf-theoretic picture
of polysemanticity. These surfaces are the manifolds for which I have explicit
cell complexes, boundary maps, cohomology computations, and
representation-theoretic decompositions already in hand.

The important caveat is that optimising for symmetry---which is what makes these
surfaces tractable---is \emph{not} the same as optimising for the quantities
that matter in the LLM setting. The Laplace--Beltrami training objective
proposed by Poole~\cite{poole2026}---minimising holonomy to make multi-agent
representations locally and globally flat---optimises for \emph{geometric}
flatness, not algebraic structure. Symmetry constrains the eigenspaces;
flatness constrains the eigenvalues. As we shall see, these pull in opposite
directions (Section~6). Nevertheless, highly symmetric surfaces make good
starting examples precisely because they are the ``spherical cows'' of
Riemannian spectral theory: simplified enough to compute with, structured enough
to reveal the essential trade-offs.

The results suggest that:

\begin{itemize}[leftmargin=*]
  \item The spectral gap $\lambda_1$ controls robustness: how resistant the
    surface is to low-frequency perturbations (analogous to prompt sensitivity
    in an LLM).
  \item Eigenspace multiplicities, forced by irreducible representations of the
    symmetry group, determine feature granularity (analogous to the resolution
    of SAE features).
  \item $H^1 = \mathbb{R}^{2g}$ is a topological invariant: the number of
    ``polysemantic handles'' that no amount of training can flatten---and, in the
    quantum coding setting, the number of logical qubits protected by the code.
  \item Maximising symmetry (Hurwitz surfaces) and maximising spectral gap
    (robustness) are in tension---a trade-off between interpretability and
    stability that may have a direct analogue in LLM design.
\end{itemize}

\section{Triangle Groups and Hyperbolic Surfaces}

A triangle group $\Delta(p,q,r) = \langle a, b \mid a^p = b^q = (ab)^r = 1 \rangle$
with $1/p + 1/q + 1/r < 1$ acts on the hyperbolic plane $\mathbb{H}^2$ by
isometries, tiling it with hyperbolic triangles of angles $\pi/p$, $\pi/q$,
$\pi/r$. If $G$ is a finite group admitting a surjection
$\Delta(p,q,r) \twoheadrightarrow G$ with torsion-free kernel
$\Gamma \trianglelefteq \Delta$, then the quotient
\[
  S = \mathbb{H}^2 / \Gamma
\]
is a compact Riemann surface with $\mathrm{Aut}(S) \supseteq G$ and genus
determined by the Riemann--Hurwitz formula:
\[
  |G| = \frac{4\pi(g - 1)}{2\pi\left(1 - \frac{1}{p} - \frac{1}{q} - \frac{1}{r}\right)}
  = \frac{2(g-1)}{1 - 1/p - 1/q - 1/r}.
\]

The \emph{Hurwitz bound} states that $|\mathrm{Aut}(S)| \leq 84(g-1)$ for
genus $g \geq 2$, achieved precisely when the triangle group is $\Delta(2,3,7)$
(since $1 - 1/2 - 1/3 - 1/7 = 1/42$ gives $|G| = 84(g-1)$). Surfaces
achieving this bound are called \emph{Hurwitz surfaces}.

\subsection{The family}

We study three surfaces, listed with their triangle group data:

\medskip
\begin{center}
\renewcommand{\arraystretch}{1.3}
\begin{tabular}{@{}lccccl@{}}
\toprule
\textbf{Surface} & \textbf{Genus} & $\Delta(p,q,r)$ & $|\mathrm{Aut}|$
  & \textbf{Hurwitz?} & \textbf{Group} \\
\midrule
Bolza & 2 & $(2,3,8)$ & 48 & No (max is 84) & $\mathrm{GL}(2, \mathbb{F}_3)$ \\
Klein quartic & 3 & $(2,3,7)$ & 168 & Yes ($= 84 \cdot 2$) & $\mathrm{PSL}(2,7)$ \\
Macbeath & 7 & $(2,3,7)$ & 504 & Yes ($= 84 \cdot 6$) & $\mathrm{PSL}(2,8)$ \\
\bottomrule
\end{tabular}
\end{center}

\medskip
The Bolza surface is the most symmetric genus-2 surface but falls short of the
Hurwitz bound: $48 < 84$. The Klein quartic is the smallest Hurwitz surface;
the Macbeath surface is the next. Each carries a regular tessellation inherited
from the triangle tiling of $\mathbb{H}^2$: the Bolza surface has a
$\{3,8\}$ triangulation (16 triangles, 6 vertices of valence~8) dual to a
$\{8,3\}$ octagonal tiling, while the Klein quartic has a $\{3,7\}$
triangulation dual to the famous $\{7,3\}$ heptagonal tiling.

\section{The Laplace--Beltrami Operator}

Equip each surface $S$ with its constant-curvature $-1$ hyperbolic metric. The
Laplace--Beltrami operator $\Delta_g$ is defined in local coordinates by
\[
  \Delta_g f = \frac{1}{\sqrt{|g|}} \sum_{i,j}
    \partial_i \left( \sqrt{|g|}\, g^{ij}\, \partial_j f \right),
\]
where $g_{ij}$ is the metric tensor and $g^{ij}$ its inverse. This generalises
the flat Laplacian $\partial^2/\partial x^2 + \partial^2/\partial y^2$ to
curved surfaces. On a compact surface without boundary, $\Delta_g$ has a
discrete spectrum:
\[
  0 = \lambda_0 < \lambda_1 \leq \lambda_2 \leq \lambda_3 \leq \cdots \to \infty.
\]

The eigenfunctions $\varphi_n$ satisfying $\Delta_g \varphi_n = -\lambda_n \varphi_n$
are the natural vibration modes of the surface---the hyperbolic analogues of
$\sin(nx)$ and $\cos(nx)$ on a circle.

\subsection{What the spectrum encodes}

\textbf{The spectral gap $\lambda_1$.} The first nontrivial eigenvalue controls
how quickly local perturbations dissipate across the surface. A large spectral
gap means the surface is ``well-mixed''---perturbations cannot remain localised.
A small gap means the surface has near-disconnected regions.

\textbf{The area.} Fixed by Gauss--Bonnet: $\mathrm{Area}(S) = 4\pi(g-1)$.
A genus-2 surface has area $4\pi$; genus~7 has area $24\pi$.

\textbf{The Yang--Yau bound.} An upper bound on $\lambda_1$:
\[
  \lambda_1 \leq \frac{8\pi(\lfloor g/2 \rfloor + 1)}{\mathrm{Area}(S)}
  = \frac{2(\lfloor g/2 \rfloor + 1)}{g - 1}.
\]
For genus~2: $\lambda_1 \leq 4$. For genus~3: $\lambda_1 \leq 8/3 \approx 2.67$.
For genus~7: $\lambda_1 \leq 4/3 \approx 1.33$.

\textbf{The Selberg bound.} For arithmetic hyperbolic surfaces (quotients by
arithmetic lattices), Selberg conjectured $\lambda_1 \geq 1/4$. The current best
result is $\lambda_1 \geq 975/4096 \approx 0.238$ (Kim--Sarnak). For random
hyperbolic surfaces of large genus, $\lambda_1 \to 1/4$ almost surely
(Brooks--Makover). So $1/4$ is the ``generic'' spectral gap.

\section{The Bolza Surface}

The Bolza surface $B$ is defined algebraically by $y^2 = x^5 - x$, with
automorphism group $\mathrm{GL}(2, \mathbb{F}_3)$ of order~48. Its fundamental
domain is a regular hyperbolic octagon with vertex angles $\pi/4$ (so that
$8 \times \pi/4 = 2\pi$ at each vertex). The octagon's sides are identified
in pairs according to the standard genus-2 surface word $aba^{-1}b^{-1}cdc^{-1}d^{-1}$.

\subsection{The spectrum}

The first several eigenvalues, computed numerically by
Strohmaier--Uski~\cite{strohmaier2013} and Cook~\cite{cook2018}, are:

\medskip
\begin{center}
\renewcommand{\arraystretch}{1.3}
\begin{tabular}{@{}ccl@{}}
\toprule
$\lambda_n$ (approx) & Multiplicity & Eigenspace structure \\
\midrule
0 & 1 & Trivial representation \\
3.839 & 2 & 2-dimensional irrep of $\mathrm{GL}(2, \mathbb{F}_3)$ \\
5.354 & 3 & 3-dimensional irrep \\
8.250 & 3 & 3-dimensional irrep \\
14.72 & 2 & 2-dimensional irrep \\
\bottomrule
\end{tabular}
\end{center}

\medskip
The spectral gap $\lambda_1 \approx 3.84$ is conjectured to be maximal among
all genus-2 hyperbolic surfaces. The Yang--Yau bound gives $\lambda_1 \leq 4$,
so the Bolza surface achieves 96\% of the theoretical maximum. This is far
above the Selberg baseline of $1/4$: the Bolza surface is a spectral outlier,
extraordinarily well-mixed for its genus.

\subsection{Representation-theoretic structure}

The group $\mathrm{GL}(2, \mathbb{F}_3)$ has order $48 = (9-1)(9-3)$. Its
irreducible representations over $\mathbb{R}$ include representations of
dimensions 1, 2, 3, and higher. The eigenspaces of $\Delta_g$ must be
$\mathrm{GL}(2, \mathbb{F}_3)$-representations, so they decompose into
direct sums of irreps.

The first nontrivial eigenspace at $\lambda_1 \approx 3.84$ is 2-dimensional:
the smallest nontrivial irrep. This means you cannot have a single
eigenfunction at $\lambda_1$---the symmetry group acts on it and generates
a 2-dimensional orbit. The symmetry forces eigenfunctions into multiplets.

\subsection{Cohomology and the $(2,3,8)$ triangle group}

The Bolza surface arises from the exact sequence
$1 \to \Gamma \to \Delta(2,3,8) \to \mathrm{GL}(2,\mathbb{F}_3) \to 1$
where $\Gamma$ is a torsion-free Fuchsian group (the fundamental group
$\pi_1(B)$). The cell complex inherited from the $(2,3,8)$ tiling gives:

\begin{itemize}[leftmargin=*]
  \item $\{3,8\}$ triangulation: 16 triangles, 24 edges, 6 vertices (valence 8)
  \item $\{8,3\}$ dual: 6 octagons, 24 edges, 16 vertices (valence 3)---the
    dual graph is the M\"{o}bius--Kantor graph
  \item $\chi = 6 - 24 + 16 = -2$, confirming genus~2
\end{itemize}

The cellular homology gives Betti numbers $\beta_0 = 1$, $\beta_1 = 4 = 2g$,
$\beta_2 = 1$, confirming $H^1(B, \mathbb{R}) \cong \mathbb{R}^4$.
The $\mathrm{GL}(2,\mathbb{F}_3)$-action on $H^1$ decomposes as
$V_2 \oplus \overline{V}_2$: a conjugate pair of 2-dimensional irreps
(see \texttt{bolza-surface/bolza\_cohomology.py} for the computation).

\section{The Klein Quartic}

The Klein quartic is the genus-3 Riemann surface defined by
$x^3 y + y^3 z + z^3 x = 0$ in $\mathbb{P}^2$, with automorphism group
$\mathrm{PSL}(2,7)$ of order~168. It is the smallest Hurwitz surface:
$168 = 84(3-1)$.

\subsection{The $(2,3,7)$ triangle group}

The Klein quartic arises from
$1 \to \Gamma \to \Delta(2,3,7) \to \mathrm{PSL}(2,7) \to 1$.
The triangle tiling of $\mathbb{H}^2$ by $(2,3,7)$ triangles descends to:

\begin{itemize}[leftmargin=*]
  \item $\{3,7\}$ triangulation: 56 triangles, 84 edges, 24 vertices (valence 7)
  \item $\{7,3\}$ dual: 24 heptagons, 84 edges, 56 vertices (valence 3)
  \item $\chi = 24 - 84 + 56 = -4$, confirming genus~3
\end{itemize}

This is the famous heptagonal tiling depicted in countless illustrations of
hyperbolic geometry. The 24 heptagons tile the genus-3 surface with 3 meeting
at each vertex.

\subsection{Representation theory of $\mathrm{PSL}(2,7)$}

$\mathrm{PSL}(2,7)$ has 6 conjugacy classes and irreducible representations of
dimensions $1, 3, \overline{3}, 6, 7, 8$. The cohomology decomposes as
\[
  H^1(\text{Klein}, \mathbb{R}) \cong \mathbb{R}^6 \cong V_3 \oplus \overline{V}_3,
\]
a conjugate pair of 3-dimensional irreps
(see \texttt{klein-quartic/klein\_quartic\_cohomology.py}). This is the
representation-theoretic signature of a Hurwitz surface: the cohomology is forced
into the smallest nontrivial irreps of the automorphism group.

\subsection{Spectral consequences of Hurwitz symmetry}

Because $\mathrm{PSL}(2,7)$ has no 2-dimensional irreps, the first nontrivial
eigenspace of $\Delta_g$ on the Klein quartic must be at least
3-dimensional---the smallest nontrivial irrep available. This is a direct
consequence of the Hurwitz bound: maximum symmetry forces the eigenspaces into
\emph{larger} multiplets.

Combined with the Yang--Yau bound $\lambda_1 \leq 8/3$ for genus~3,
this constrains the spectrum tightly. The Klein quartic cannot have a spectral
gap as large as the Bolza surface ($\lambda_1 \leq 2.67$ vs.\ the Bolza's
$\lambda_1 \approx 3.84$), but its eigenspaces are more highly structured.

\section{The Macbeath Surface}

The Macbeath surface is the genus-7 Hurwitz surface with automorphism group
$\mathrm{PSL}(2,8)$ of order $504 = 84 \cdot 6$. It arises from the same
$(2,3,7)$ triangle group as the Klein quartic, but with a different
(larger) quotient.

\subsection{Representation theory of $\mathrm{PSL}(2,8)$}

$\mathrm{PSL}(2,8)$ has 9 conjugacy classes with irreducible representations of
dimensions $1, 7, 7, 7, 8, 8, 8, 9, 9$. Notably, the smallest nontrivial
irrep is 7-dimensional. This forces each nontrivial eigenspace of $\Delta_g$ to
have dimension at least~7.

The cohomology is $H^1 \cong \mathbb{R}^{14}$, and it decomposes into
$\mathrm{PSL}(2,8)$-irreps. The 14-dimensional $H^1$ provides 14 independent
topological obstructions---14 independent directions in which local frames
cannot be made globally compatible.

\subsection{The spectral gap shrinks}

The Yang--Yau bound for genus~7 is $\lambda_1 \leq 4/3 \approx 1.33$. Combined
with eigenspaces of dimension $\geq 7$, the Macbeath surface has a coarser
spectral resolution than either the Bolza or Klein surfaces: fewer distinct
eigenvalues in any given range, each with high multiplicity.

\section{The Hurwitz Paradox: Symmetry vs.\ Spectral Gap}

\subsection{The tension}

Comparing the three surfaces reveals a fundamental tension:

\medskip
\begin{center}
\renewcommand{\arraystretch}{1.3}
\begin{tabular}{@{}lccccc@{}}
\toprule
\textbf{Surface} & $g$ & $|\mathrm{Aut}|$ & Hurwitz?
  & $\lambda_1$ bound & Min eigenspace dim \\
\midrule
Bolza & 2 & 48 & No & $\leq 4$ & 2 \\
Klein & 3 & 168 & Yes & $\leq 2.67$ & 3 \\
Macbeath & 7 & 504 & Yes & $\leq 1.33$ & 7 \\
\bottomrule
\end{tabular}
\end{center}

\medskip
More symmetry means:
\begin{itemize}[leftmargin=*]
  \item Larger minimum eigenspace dimensions (eigenspaces must be full irreps)
  \item A lower upper bound on $\lambda_1$ (via Yang--Yau, since genus grows with symmetry)
  \item Coarser spectral resolution (fewer distinct eigenvalues, each highly degenerate)
\end{itemize}

Less symmetry means:
\begin{itemize}[leftmargin=*]
  \item Smaller eigenspace dimensions (finer-grained modes)
  \item A potentially larger spectral gap (the Bolza surface nearly saturates Yang--Yau)
  \item More distinct eigenvalues in any given spectral range
\end{itemize}

\subsection{Symmetry maximises structure; spectral gap maximises mixing}

The Hurwitz bound $84(g-1)$ maximises the number of symmetries a surface of
genus $g$ can carry. But maximising symmetry is not the same as maximising the
spectral gap. The Bolza surface has less symmetry than the Hurwitz bound
allows, but it has the largest spectral gap of any genus-2 surface.

The two criteria optimise different things:
\begin{itemize}[leftmargin=*]
  \item \textbf{Symmetry (Hurwitz)} $\to$ the eigenspaces are maximally
    constrained by representation theory $\to$ the spectral structure is most
    ``interpretable'' (fewest free parameters, everything determined by group theory).
  \item \textbf{Spectral gap} $\to$ perturbations dissipate maximally fast
    $\to$ the surface is most ``robust'' (hardest to excite nontrivial modes).
\end{itemize}

These are in tension because achieving the Hurwitz bound requires genus to grow
(for larger groups), which pushes the Yang--Yau bound down and forces the
spectral gap to shrink.

\section{Topological Constraints: Why Curvature Cannot Be Eliminated}

\subsection{Gauss--Bonnet}

On any compact surface of genus $g \geq 2$ with constant curvature $K = -1$:
\[
  \int_S K\, dA = 2\pi\chi(S) = 4\pi(1 - g).
\]
The total curvature is a topological invariant. You can redistribute curvature
(make some regions flatter and others more curved), but you cannot eliminate it.
The negative curvature is topologically forced.

\subsection{$H^1$ and the handles}

The first cohomology $H^1(S, \mathbb{R}) \cong \mathbb{R}^{2g}$ records the
topological obstructions. Each handle of the surface contributes two
independent directions to $H^1$, and each direction represents a way in which
parallel transport around a non-contractible loop can fail to return a vector
to its starting orientation.

\medskip
\begin{center}
\renewcommand{\arraystretch}{1.3}
\begin{tabular}{@{}lccl@{}}
\toprule
\textbf{Surface} & Genus & $\dim H^1$ & $H^1$ as Aut-representation \\
\midrule
Bolza & 2 & 4 & $V_2 \oplus \overline{V}_2$ under $\mathrm{GL}(2,\mathbb{F}_3)$ \\
Klein & 3 & 6 & $V_3 \oplus \overline{V}_3$ under $\mathrm{PSL}(2,7)$ \\
Macbeath & 7 & 14 & Irreps of $\mathrm{PSL}(2,8)$ \\
\bottomrule
\end{tabular}
\end{center}

\medskip
The pattern is clear: $H^1$ always decomposes into conjugate pairs of the
\emph{smallest} nontrivial irreps of the automorphism group. For the Bolza
surface, the smallest is 2-dimensional; for the Klein quartic,
3-dimensional; for Macbeath, the 7-dimensional irreps. The topology forces $H^1$
to fill out complete irreducible representations---it cannot be ``half'' of an
irrep.

\subsection{Holonomy as parallel transport failure}

The holonomy group of a Riemannian manifold measures what happens when you
parallel-transport a tangent vector around a closed loop. On a flat manifold,
the vector returns unchanged (trivial holonomy). On a curved manifold, it
returns rotated.

On a genus-$g$ surface with constant curvature $-1$, the holonomy around the
$2g$ generators of $\pi_1(S)$ is nontrivial. These are the ``handles'' of
the surface, and the holonomy around each handle is determined by the
hyperbolic translation length of the corresponding element of the Fuchsian
group $\Gamma$.

Trivial holonomy would require the surface to be flat---but Gauss--Bonnet
forbids this for $g \geq 2$. The holonomy is therefore topologically forced:
a consequence of the genus, not a failure of optimisation.

\section{Eigenfunctions as Irreducible Representations}

\subsection{The general principle}

If a group $G$ acts on a Riemannian manifold $S$ by isometries, then $G$ acts on
each eigenspace $E_\lambda = \{f : \Delta_g f = -\lambda f\}$. By Schur's lemma,
$E_\lambda$ decomposes into a direct sum of irreducible $G$-representations:
\[
  E_\lambda \cong \bigoplus_i V_{\rho_i}^{\oplus m_i}.
\]
The multiplicities $m_i$ are the Fourier coefficients of the spectral
decomposition with respect to the representation ring of $G$.

This means:
\begin{itemize}[leftmargin=*]
  \item Each eigenvalue ``carries'' one or more irreps of the symmetry group.
  \item The minimum eigenspace dimension equals the dimension of the smallest
    irrep that appears at that eigenvalue.
  \item Two eigenvalues can only be equal if their eigenspaces merge into a
    single (larger) $G$-representation.
\end{itemize}

\subsection{Consequences for spectral resolution}

A group with large minimum irrep dimension forces coarse spectral resolution.
$\mathrm{PSL}(2,8)$ has minimum nontrivial irrep dimension~7, so the Macbeath
surface cannot have an eigenspace smaller than 7-dimensional (apart from
the trivial eigenvalue $\lambda_0 = 0$). In contrast,
$\mathrm{GL}(2,\mathbb{F}_3)$ has a 2-dimensional irrep, so the Bolza surface
can have eigenspaces as small as 2-dimensional.

This is a direct trade-off: more symmetry buys you more algebraic structure
(the eigenspaces are determined by representation theory), but at the cost of
spectral resolution (each eigenvalue accounts for more dimensions). The
``finest-grained'' spectral analysis is available on the least symmetric surfaces.

\section{Conclusion}

The Laplace--Beltrami operator on a compact hyperbolic surface is the nexus
where topology, geometry, and representation theory meet:

\medskip
\begin{center}
\renewcommand{\arraystretch}{1.3}
\begin{tabular}{@{}ll@{}}
\toprule
\textbf{Property} & \textbf{What it determines} \\
\midrule
$\lambda_1$ (spectral gap) & Mixing rate; robustness to perturbation \\
Eigenspace dimensions & Forced by irreps of $\mathrm{Aut}(S)$ \\
$H^1 \cong \mathbb{R}^{2g}$ & Topological obstructions (handles) \\
$H^1$ as $\mathrm{Aut}$-rep & How symmetry organises the obstructions \\
Yang--Yau bound & Maximum possible spectral gap for given genus \\
Selberg bound ($1/4$) & Generic large-genus baseline \\
Holonomy & Parallel transport failure around handles \\
\bottomrule
\end{tabular}
\end{center}

\medskip
The Hurwitz paradox---that maximising symmetry and maximising spectral gap are
in tension---reveals that these surfaces optimise different criteria. Hurwitz
surfaces maximise algebraic structure (the representation theory determines
everything). Non-Hurwitz surfaces like the Bolza can maximise spectral gap
(robustness to perturbation) at the cost of less algebraic constraint.

The topological invariant $H^1 = \mathbb{R}^{2g}$, decomposed into irreducible
representations of the automorphism group, is the fundamental datum: it measures
what cannot be changed by any geometric deformation, only by changing the
topology itself. On a genus-$g$ surface, there are exactly $2g$ independent
directions in which curvature is topologically forced and holonomy is necessarily
nontrivial. This is a theorem, not an approximation.

\begin{thebibliography}{9}
\bibitem{langer2026llm}
R.~Langer.
\emph{LLMs as Resonant Structures: From Green's Functions to Nonlinear
Spectral Theory}. 2026.

\bibitem{strohmaier2013}
A.~Strohmaier and V.~Uski.
An algorithm for the computation of eigenvalues, spectral zeta functions
and zeta-determinants on hyperbolic surfaces.
\emph{Communications in Mathematical Physics}, 317(3):827--869, 2013.

\bibitem{cook2018}
J.~Cook.
Properties of eigenvalues on Riemann surfaces with large symmetry groups.
PhD thesis, Loughborough University, 2018.

\bibitem{kac1966}
M.~Kac.
Can one hear the shape of a drum?
\emph{The American Mathematical Monthly}, 73(4):1--23, 1966.

\bibitem{hurwitz1893}
A.~Hurwitz.
\"{U}ber algebraische Gebilde mit eindeutigen Transformationen in sich.
\emph{Mathematische Annalen}, 41:403--442, 1893.

\bibitem{yang_yau1980}
P.~C.~Yang and S.-T.~Yau.
Eigenvalues of the Laplacian of compact Riemann surfaces and minimal submanifolds.
\emph{Annali della Scuola Normale Superiore di Pisa}, 7(1):55--63, 1980.

\bibitem{selberg1965}
A.~Selberg.
On the estimation of Fourier coefficients of modular forms.
\emph{Proc.\ Symp.\ Pure Math.}, 8:1--15, 1965.

\bibitem{brooks_makover2004}
R.~Brooks and E.~Makover.
Random construction of Riemann surfaces.
\emph{Journal of Differential Geometry}, 68(1):121--157, 2004.

\bibitem{langer2026css}
R.~Langer.
\emph{CSS Quantum Codes from Surface Tilings: A Lean~4 Formalisation}.
2026.

\bibitem{poole2026}
A.~Poole.
Personal communication, August 2026.
\end{thebibliography}

\end{document}
